\documentclass[10pt,a4paper]{amsart}
\usepackage[margin=19mm]{geometry}
\usepackage{amsmath,amssymb,mathtools}
\usepackage[scr=rsfs,cal=euler]{mathalfa}
\usepackage{xfrac}
\usepackage[unicode,hidelinks,bookmarks=false]{hyperref}
\newcommand{\ie}{i.e.\ }
\newcommand{\cf}{cf.\ }
\newcommand{\resp}{respectively\ }
\newcommand{\Subsection}{\subsection}
\providecommand{\nobreakdash}{\nobreak\hbox{-}}
\setlength{\emergencystretch}{2em}
\multlinegap0pt
\usepackage{dsfont}
\usepackage{amssymb}
\usepackage{xcolor}
\usepackage[shortlabels]{enumitem}
\newtheorem{thm}{Theorem}
\newcommand{\C}{\mathbb{C}}
\newcommand{\Cn}{\mathbb{C}^n}
\title{A Canonical Characterization of Normal Functions}
\author{Peter V Dovbush, Steven G Krantz}
\date{Source: arXiv:2601.20603v2 (CC BY 4.0)}
\begin{document}
\maketitle

\noindent\emph{Source and licence.} A Canonical Characterization of Normal Functions. Version arXiv:2601.20603v2 (CC BY 4.0), distributed by arXiv under the Creative Commons Attribution 4.0 International licence, http://creativecommons.org/licenses/by/4.0/, as recorded in the arXiv licence metadata for this exact version and shown on the abstract page https://arxiv.org/abs/2601.20603v2. Full licence notice: \texttt{licenses/arxiv-2601.20603-CC-BY-4.0.txt}.\par\smallskip

\noindent\textit{Editorial scope.} Counted unit: the article's thm thm:norm-manif-1 (source0.tex lines 269--277). The article's complete original proof of it is delivered with it (source0.tex lines 279--342, one \texttt{proof} environment of 64 lines). No mathematics is written, repaired or re-derived: the statement and the proof are the article's own lines, in the article's own notation, and no retained line is substituted or re-flowed.  Retained uncounted as the source-local context the proof needs: lines 212--220.  Nothing was omitted from any retained span, so the package carries no omission record.  No retained line contains a citation, so the wrapper carries no bibliography and no bibliography entry.  No retained line contains a figure environment, an image-inclusion command or drawing code, so no image is delivered and the package declares no asset dependency.  Editorial formatting only, in addition: an AMS \texttt{amsart} preamble that restates the article's own declarations and macros used by the retained lines, namely the environments \texttt{thm} on separate counters, together with the standard packages those lines need.  The retained environments print in this document as Theorem 1 in order of appearance, whereas the article prints the same items with the numbering of its own full document.  The complete native source of the article as distributed in the arXiv e-print for this exact version is bundled unchanged as \texttt{sources/193538/source0.tex}.
    \begin{thm}
        \label{thm:marty:1}
        A family $\mathcal{F}$ of functions holomorphic on $\Omega$ is normal on $\Omega\subset \Cn$ if and only if for each compact subset $K\subset \Omega$ there exists a constant $M(K)$ such that at each point $z\in K$
        \begin{equation}
            f^\sharp(z)\leq M(K)
            \label{eq:marty:1}
        \end{equation}
        for all $f\in \mathcal{F}$.
    \end{thm}

    \begin{thm}
        \label{thm:norm-manif-1}
        If a subset $\mathcal{F}$ of $\mathcal{O}(\mathbb{B}, \C)$ is an $\operatorname{Aut}(\mathbb{B})$-invariant normal family, then there exists a constant $C$ such that
        \begin{equation}
            L_z(\log(1+|f|^2),v)\leq C \cdot ds^2(z,v)
            \label{eq:norm-inequality}
        \end{equation}
        for every $f \in \mathcal{F}$.
    \end{thm}

    \begin{proof}
        For $z\in\mathbb{B}$, define
        \[
            C(z) := \sup_{\substack{f\in\mathcal{F}\\v\ne0}}\frac{ L_z\left(\log(1+|f|^2),v\right) }{ ds^2(z,v) }\,.
        \]

        We first show that $C(0)<\infty$.

        Since $\mathcal{F}$ is a normal family, Marty's theorem (Theorem~\ref{thm:marty:1}) implies that the spherical derivatives of the functions in $\mathcal{F}$ are locally uniformly bounded.
        In particular, there exists a constant $M<\infty$ such that
        \[
            L_0\left(\log(1+|f|^2),v\right)\leq M|v|^2
        \]
        for every $f\in\mathcal{F}$ and every $v\in\C^n$.

        We have
        \[
            2|v|^2=ds^2(0,v)
        \]
        for every $v\in\C^n$.
        Consequently,
        \[
            L_0\left(\log(1+|f|^2),v\right) \leq M|v|^2\leq Mds^2(0,v)\,.
        \]
        Hence
        \[
            C(0)\leq M<\infty\,.
        \]

        We next show that $C$ is invariant under $\operatorname{Aut}(\mathbb{B})$.
        Let $g\in\operatorname{Aut}(\mathbb{B})$.
        Since $\mathcal{F}$ is $\operatorname{Aut}(\mathbb{B})$-invariant, $f\circ g\in\mathcal{F}$ whenever $f\in\mathcal{F}$.
        By the chain rule,
        \[
            L_z\left(\log(1+|f\circ g|^2),v\right) = L_{g(z)}\left(\log(1+|f|^2),g_{*}v\right)\,.
        \]
        By the definition of $C(g(z))$,
        \[
            L_{g(z)}\left(\log(1+|f|^2),g_{*}v\right) \leq C(g(z))\,ds^2(g(z),g_{*}v)\,.
        \]
        The $\operatorname{Aut}(\mathbb{B})$-invariance of the Bergman metric gives
        \[
            ds^2(g(z),g_{*}v)=ds^2(z,v)\,.
        \]
        Therefore
        \[
            L_z\left(\log(1+|f\circ g|^2),v\right) \leq C(g(z))\,ds^2(z,v)\,.
        \]
        Taking the supremum over $f\in\mathcal{F}$ and $v\neq0$ gives
        \[
            C(z)\leq C(g(z))\,.
        \]
        Replacing $g$ by $g^{-1}$ gives the reverse inequality.
        Hence
        \[
            C(z)=C(g(z))\,.
        \]
        Since $\mathbb{B}$ is homogeneous under $\operatorname{Aut}(\mathbb{B})$, $C(z)$ is constant on $\mathbb{B}$.
        Since $C(0)<\infty$, this constant is finite, and therefore
        \[
            L_z\left(\log(1+|f|^2),v\right) \leq C\,ds^2(z,v)
        \]
        for every $f\in\mathcal{F}$, every $z\in\mathbb{B}$, and every $v\in\C^n$.
    \end{proof}

\end{document}
